\section{The Light of the Imperium}

\textit{In hoc signo vinces.}

\begin{quotex}
Europe is haunted by the shadow of the Emperor. One senses his absence just as vividly as in former times one sensed his presence. Because the emptiness of the wound \emph{speaks}, what we miss knows how to make us sense it.
\end{quotex}

Last month Pope Benedict visited England\footnote{\url{https://www.ft.com/cms/s/2/0e6f301c-baea-11df-9e1d-00144feab49a,dwp\_uuid=a712eb94-dc2b-11da-890d-0000779e2340.html}} to reclaim Britain for the Roman Empire. The particular circumstances are irrelevant. What would have been unthinkable not so long ago, became possible. The ultimate path of Britain was set the day of its schism. Its monarchs no longer have authority, no one accepts Prince Charles as a possible king, its local ``church'' is rootless, and the once great nation is annihilating itself.

Life, human life, cannot grow and thrive in the shadows. As we prepare to remember the dead, we will let you meditate on the thoughts of those who still can recall the light of the Imperium.

\begin{quotationx}
The post of the Emperor does not belong to those who desire it or to the choice of the people, it is reserved to the choice of heaven alone. It has become occult.

The ruler of men would no longer be a man, but a being of higher level, as a ``man'' --- even if, on the exterior, he maintains, more or less, a common human appearance --- through the fact that the hierarchy, whose members are at this point consciousnesses, is immaterial and cannot be distinguished by any physically visible feature. As such, the ruler would no longer be compared, for instance, to a hand which wants to make itself master over the whole body, but should rather be compared to the organic unity of the body which, in a higher, incorporeal synthesis, comprises the hand and everything else.

The \emph{post} of the Emperor ... what an abundance of ideas concerning the post --- its historical mission, its functions in the light of natural right, and its role in the light of divine fight --- of the Emperor of Christendom are to be found among mediaeval authors!

The idea of ``Lord of humanity'' is by no means one invented by us: it corresponds precisely to the primordial Aryan concept of the cakravarti Lord of the World , which, in the symbolic terms of saga and myth, was constantly connected with the real or legendary figures of great rulers, from Alexander the Great\footnote{\url{http://www.gornahoor.net/?p=331}} to King Arthur\footnote{\url{http://www.gornahoor.net/?p=277}} and Emperor Frederick II\footnote{\url{http://www.gornahoor.net/?p=826}}.

Authors of the Middle Ages could not imagine Christendom without an Emperor. Because if the world is governed hierarchically, the \emph{Sanctum Imperium} cannot be otherwise. Hierarchy is a pyramid which exists only when it is complete. And it is the Emperor who is at its summit. Then comes the kings, dukes, noblemen, citizens and peasants.

We are intransigent affirmers of the necessity of hierarchy, we maintain that this hierarchy must be built dynamically and freely, through natural relationships of individual intensity. Primitive aristocracies formed like this --- even where a supernatural principle did not impose them directly --- not by election and recognition from below, but by the direct self-assertion of individuals capable of a degree of resistance, of responsibility, of a heroic, generous, full, and dangerous life, which the others were not capable of.

With that, keep in mind that we do not want to abolish ``man'', that is, that consciousness of freedom, individuality, and autonomy of individuals gained against the primitive, indistinct, mediumistic sociality. A true King never desires shadows, puppets, and automatons as subjects, but rather he desires individuals, warriors, living, and strong beings; and in fact, his pride would be to feel himself to be a King of kings.

The Emperor reigns by pure authority; he reigns over free beings, not by means of the sword, but by means of the sceptre. Just as the world is ruled by the cross, so the power of the Emperor over the terrestrial globe is subject to the sign of the cross.

Hierarchy can not be restored by violence, the control of needs, or the interplay of passions, interests and ambitions, but only by the free and spontaneous recognition which springs from the sense of values and of transcendent forces, from faithfulness toward one's own way of being, whatever it might be, from consciousness of nature, dignity and quality. An organic, direct, real, hierarchy: freer and harder than any other.
\end{quotationx}

All quotations are take from \textbf{Meditations on the Tarot} by Valentin Tomberg and from \textbf{Pagan Imperialism} by Julius Evola.


\flrightit{Posted on 2010-10-30 by Cologero}

\begin{center}* * *\end{center}

\begin{footnotesize}
\begin{sffamily}
\texttt{EXIT on 2010-10-31 at 10:00 said:}

``reclaim Britain for the Roman Empire''

Don't you mean for the neo-Marxist Catholic church?

``The ultimate path of Britain was set the day of its schism.''

Schisms between Catholic and Protestant or Anglican are like debates between the mainstream political parties: there's not an ounce of difference between them; they're both corrupt, irrational and illegitimate. The Christian church sealed its fate in the early schism between Roman Catholicism and the Alexandrian Gnostics. The Catholic church suffers from the same problem that modern royalty faces and that is the total lack of initiation and spiritual realization which alone makes it lawful. Catholicism and the regality have become exoterisms unto them selves completely divorced from true metaphysics. I doubt anything will ever change this fact, but surely not the bunch of PR stunts which the church is reduced to.

``Hierarchy can not be restored by violence''

I would have to disagree with that statement. People today are predisposed to ignorance on account of their genes deciding behavior. Hence, there won't be any ``free and spontaneous recognition which springs from the sense of values and of transcendent forces''; there will instead be more of the same infantile garbage dictated by religious fundamentalist wannabes, repulsive moralists, and those who think turning the other cheek and giving up the shirt off your back to thieves is spirituality.

\hfill

\texttt{Cologero on 2010-10-31 at 11:27 said:}

Of course, we were speaking figuratively and referred to the Pope's visit solely as an archetype repeating in history. As such, it is a point of departure to go back to understand the archetypes themselves, viz., the posts of Emperor and Pope. If the post of Emperor is vacant, there is nothing to go back to at the moment. Similarly, the post of Pope may also be vacant, since as you pointed out, it no longer reflects a Traditional spirituality. It is exceedingly difficult to lead readers to abstract thinking which means the ability to think in terms of ideals and principles rather that in terms of particular contingent, historical, and empirical facts. Therefore, we are interested in the thoughts and ideas of the greatest minds, not ``infantile garbage'' nor ``fundamentalist wannabes''.

The moral questions addressed are the notions of legitimacy of authority and fidelity. How anyone moves from there to a discussion of cheeks and shirts is disconcerting and somewhat annoying. Perhaps we have been too subtle.

First of all, the position of Gornahoor is that the Middle Ages represents a Traditional civilization, a position also held by Evola and Guenon. So a one-liner about Alexandrian Gnostics demonstrates nothing much. Second, we, along with Evola, adhere to Maistre's dictum that what we need is not a revolution in the opposite direction, but the opposite of a revolution. So a revolutionary call is simply another symptom of modernity, characteristic of the lower classes.

A so-called ``hierarchy'' based on violence would be a hierarchy of the most violent, not of the most noble; citizens would accept it out of fear, not faithful allegiance. So the question of the legitimate source of authority must be addressed. We assert, along with Evola, Guenon, and Tomberg, that the source is transcendent, not ``genetic''. We linked to a news article regarding freedom of conscience for a reason. Only a free man can act morally, and one of his first moral acts is fidelity to legitimate authority. Now there may be rebellion against legitimate authority; this is called revolution. But in the case of the very lack of legitimate authority, there is no requirement for fidelity, yet revolution will accomplish nothing. There must be a complete change of mind, or wordview --- ``the subtle rules the dense''. The investigation of this task is the impetus behind Gornahoor.

\hfill

\texttt{EXIT on 2010-10-31 at 12:40 said:}

The circumstances are not irrelevant, nor are the points which you brought up beyond reproach. It seems though you like to take my counterpoints out of context as a defense of your own. As for ``hierarchy based on violence'', I said nothing of the sort, rather it is possible to ``restore'' hierarchy by such. As for Gornahoor's ``traditional'' worldview, there is much occultism thrown in here and there, so I suspect that a complete change of worldview is difficult even for you. How much more so it would be for those whose lives revolve around machines. One shouldn't expect the Western world to spontaneously restore its own traditions. Moreover, a restoration of Catholicism would by far be negative as it has never valued the free individual but desires mindless slaves to an irrational morality; and this is the reason that the church came into conflict with true regality.

``The moral questions addressed are the notions of legitimacy of authority and fidelity. How anyone moves from there to a discussion of cheeks and shirts is disconcerting and somewhat annoying. Perhaps we have been too subtle.''

Or perhaps you let your intellectual pride get in the way. Those ``cheeks and shirts'' are teachings from Jesus formulated in opposition to Pharisee rule, so pertinent to this discussion where Christianity is in question.

What is annoying is this continued assertion that Christianity can be the basis for a renewed traditional Western world without taking into consideration any of the real and practical circumstances of the time period. This isn't the Middle Ages, and moreover, even in the Middle Ages we say a ``hierarchy'' enforced almost solely through violence. But let's not gloss over these ``circumstances.''

\hfill

\texttt{EXIT on 2010-10-31 at 12:46 said:}

``even in the Middle Ages we say a hierarchy enforced almost solely through violence.''

That should have read ``saw a hierarchy... ''

\hfill

\texttt{Matt on 2010-10-31 at 17:01 said:}

EXIT,

I would say part of what is expressed on Gornahoor can be considered occult (in the sense of its traditional definition, that which is hidden) but not in the way of the ``occultism'' that you might be referring to.

As for the hierarchy in the Middle Ages, it was enforced partly through violence, but not solely or almost solely through those means.

\hfill

\texttt{EXIT on 2010-11-02 at 09:18 said:}

The question is do you know the difference?

\hfill

\texttt{kadambari on 2010-11-02 at 10:22 said:}

and those who think turning the other cheek and giving up the shirt off your back to thieves is spirituality.

Actually EXIT has a good point. This kind of irrationality was obtained by Gandhi from his Christian readings (the sources of his ideas were from a variety of sources, for example his dietary habits such as avoiding tea was obtained in a South African prison, and not from Hinduism as people like to think), as they are nowhere to be found in the native religions. In a way, I can sympathize with what he is saying.

\hfill

\texttt{James O'Meara on 2010-11-02 at 12:06 said:}

``This kind of irrationality was obtained by Gandhi from his Christian readings (the sources of his ideas were from a variety of sources, for example his dietary habits such as avoiding tea was obtained in a South African prison, and not from Hinduism as people like to think), as they are nowhere to be found in the native religions.''

As I've mentioned before, Alain Danielou has some interesting passages in his autobiography on such British-aping outcastes as Nehru and Gandhi, whom the media promoted as ``the voice of the Indian masses'' etc.

\hfill

\texttt{kadambari on 2010-11-02 at 13:28 said:}

Well I do not deny Gandhi for what he did achieve, he could have gotten scared and run away (like most Indians), but he did go back and do something, even if the results came out the way they did. One has to respect that he was an extremely determined man; but this is all. When I read his autobiography, I just could not see him as a leader, the book speaks for itself.

A bad form of governance can turn the nation into a basket case. I remember my old American friends telling me that they used to collect money to send to the Chinese back then, and it was the Chinese that had a population problem. Look at them today, sure they are on a comsumption spree and have adopted Communism, but they are focused on themselves and on attaining their potential which is quite high judging from their history.. India never quite had a revolution in the sense of a complete change of heart, after independence all the previous structures remained intact. Protesting in non-violent marches against the British by petit-bourgeois type people wearing home spun clothing does not constitute a revolution. India adopted willy nilly a form of governance from their British masters, without any thought of how this is applicable to India, and forgetting they are fast forwarding the several centuries it took for the British to obtain that form of governance. Today India is a mediocre country known for Bollywood, the cricket arena, and for people who need to be elsewhere than in India to make something out of themselves. But if you point this to them, they will think you are anti-Indian. Anyway, when a nation only copies others and cannot think for itself what its needs are and has to care for the approval of others, then it is a deracinated nation in every way. Since it is unlikely there will be any kind of revolution, what is likely is 8\% growth which is accompanied by massive levels of poverty all side by side... They might wake up one day, but the demography will be against them by then... 

\hfill

\texttt{kadambari on 2010-11-02 at 13:56 said:}

@James

After independence the right kinds of leaders were not even lacking, people like Sardar Patel, Bose, Savarkar and such, people with vision. But these people and their way of thinking was completely sidelined. Even today, there are many really intelligent people who choose to stay there, but the odds against them are very high... Since we private individuals are not Genghis Khans to change things, I guess one can just keep informed and awake and preserve the good things of one's tradition in the small way one can... 

\hfill

\texttt{Matt on 2010-11-02 at 19:59 said:}

Exit,

I'm guessing that comment was directed towards me.

There is a difference between say the metaphysical idea of the hidden centers of the body seen in Tantra, and doing spells and trying to conjure up beings (what occultists concern themselves with).

\hfill

\texttt{kadambari on 2010-11-03 at 09:37 said:}

King Arthur--

The Ossetian origins of the King Arthur legend is interesting... 

From Scythia to Camelot:
\url{http://www.fravahr.org/spip.php?article494}


\hfill

\texttt{kadambari on 2010-11-03 at 13:00 said:}

I am not a scholar in the above, but this seems like an interesting book, although I do not know enough about the topic to read it critically; cultures were much more fluid back then and there were more interactions between East and West than we know. The Buddha is also said to have Scythian origins... so it is all interesting... 

\hfill

\texttt{kadambari on 2010-11-04 at 12:23 said:}

Interesting list for Europeans. As far as magnanimity (in the old sense of megalopsuchia greatness of soul) and nobility goes, I think Ashoka and Cyrus the Great are models for Easterners in the sense of being enlightened emperors... 

\hfill

\end{sffamily}
\end{footnotesize}

